\documentclass{article}
\usepackage{amsmath, amssymb, amsthm}

\setlength{\parindent}{0pt}

\newtheorem*{definition}{Definition}
\newtheorem{claim}{Claim}

\begin{document}

\begin{definition}
    A derangement is a permutation $(x_1, x_2, \dots, x_n)$ of $\{1, 2, \dots, n\}$ such that $x_i \neq i$ for all $i$.
\end{definition}

Let $S_i$ be the set of all permutations $(x_1, x_2, \dots, x_n)$ that are not derangements because $x_i = i$. Then the 
set of non-derangements is
\[
\bigcup_{i=1}^n S_i
\]

\begin{claim}
    $|S_i|$ is $(n - 1)!$.
\end{claim}

\begin{proof}
    Since every permutation in $S_i$ satisfies $x_i = i$, it leaves the remaining $n - 1$ positions to be filled by the 
    remaining $n - 1$ elements of $\{1, 2, \dots, n\} \setminus \{i\}$, in any order. There are $(n - 1)!$ ways to 
    arrange these, so
    \[
    |S_i| = (n - 1)!
    \]

    \medskip

    Hence $|S_i|$ is $(n - 1)!$.
\end{proof}

\begin{claim}
    $|S_i \cap S_j|$ where $i \neq j$ is $(n - 2)!$.
\end{claim}

\begin{proof}
    Since every permutation in $S_i \cap S_j$ satisfies $x_i = i$ and $x_j = j$, it leaves the remaining $n - 2$ 
    positions to be filled by the remaining $n - 2$ elements of $\{1, 2, \dots, n\} \setminus \{i, j\}$, in any order. 
    There are $(n - 2)!$ ways to arrange these, so
    \[
    |S_i \cap S_j| = (n - 2)!
    \]

    \medskip

    Hence $|S_i \cap S_j|$ where $i \neq j$ is $(n - 2)!$.
\end{proof}

\begin{claim}
    $|S_{i_1} \cap S_{i_2} \cap \dots \cap S_{i_k}|$ where $i_1, i_2, \dots, i_k$ are distinct is $(n - k)!$.
\end{claim}

\begin{proof}
    Since every permutation in $S_{i_1} \cap S_{i_2} \cap \dots \cap S_{i_k}$ satisfies $x_{i_1} = i_1$, $x_{i_2} = i_2$, 
    \dots, $x_{i_k} = i_k$, it leaves the remaining $n - k$ positions to be filled by the remaining $n - k$ elements of 
    $\{1, 2, \dots, n\} \setminus \{i_1, i_2, \dots, i_k\}$, in any order. There are $(n - k)!$ ways to arrange these, so 
    \[
    |S_{i_1} \cap S_{i_2} \cap \dots \cap S_{i_k}| = (n - k)!
    \]

    \medskip

    Hence $|S_{i_1} \cap S_{i_2} \cap \dots \cap S_{i_k}|$ where $i_1, i_2, \dots, i_k$ are distinct is $(n - k)!$.
\end{proof}

\begin{claim}
    The number of non-derangements of $\{1, 2, \dots, n\}$ is
    \[
    \sum_{k=1}^{n}(-1)^{k+1}\left(\frac{n!}{k!}\right)
    \]
\end{claim}

\begin{proof}
    By the inclusion-exclusion principle,
    \[
    \left|\bigcup_{i=1}^n S_i\right| = \sum_{k=1}^n (-1)^{k+1}\sum_{1 \le i_1 < \dots < i_k \le n}
    |S_{i_1} \cap \dots \cap S_{i_k}|
    \]
    By the previous claim,
    \[
    |S_{i_1}\cap \dots \cap S_{i_k}| = (n - k)!
    \]
    where $i_1, \dots, i_k$ are distinct. Since there are $\binom{n}{k}$ choices of such indices,
    \[
    \left|\bigcup_{i=1}^n S_i\right| = \sum_{k=1}^n (-1)^{k+1}\binom{n}{k}(n - k)!
    \]
    Finally,
    \[
    \begin{aligned}
    \binom{n}{k}(n - k)! &= \frac{n!}{k!(n - k)!}(n - k)! \\
    &= \frac{n!}{k!}
    \end{aligned}
    \]
    so
    \[
    \left|\bigcup_{i=1}^n S_i\right| = \sum_{k=1}^n (-1)^{k+1}\left(\frac{n!}{k!}\right)
    \]

    \medskip

    Hence the number of non-derangements of $\{1, 2, \dots, n\}$ is
    \[
    \sum_{k=1}^{n}(-1)^{k+1}\left(\frac{n!}{k!}\right)
    \]
\end{proof}

\begin{claim}
    There are $\binom{n}{k}$ terms with the form $|S_{i_1} \cap S_{i_2} \cap \dots \cap S_{i_k}|$ in the 
    inclusion-exclusion expansion leading to the expression
    \[
    \sum_{k=1}^{n}(-1)^{k+1}\left(\frac{n!}{k!}\right)
    \]
\end{claim}

\begin{proof}
    Each such term is uniquely determined by the choice of indices $i_1, \dots, i_k$ which are all distinct. This is 
    equivalent to choosing a $k$-element subset of $\{1, 2, \dots, n\}$. There are $\binom{n}{k}$ such choices, thus 
    the number of terms is $\binom{n}{k}$.

    \medskip

    Hence there are $\binom{n}{k}$ terms with the form $|S_{i_1} \cap S_{i_2} \cap \dots \cap S_{i_k}|$ in the 
    inclusion-exclusion expansion leading to the expression
    \[
    \sum_{k=1}^{n}(-1)^{k+1}\left(\frac{n!}{k!}\right)
    \]
\end{proof}

\begin{claim}
    The number of non-derangements is
    \[
    n!\left(\frac{1}{1!} - \frac{1}{2!} + \frac{1}{3!} - \dots \pm \frac{1}{n!}\right)
    \]
    thus the number of derangements is
    \[
    n!\left(1 - \frac{1}{1!} + \frac{1}{2!} - \frac{1}{3!} + \dots \pm \frac{1}{n!}\right)
    \]
\end{claim}

\begin{proof}
    From Claim 4, the expression for the number of non-derangements can be rewritten as follows:
    \[
    \begin{aligned}
    \sum_{k=1}^{n}(-1)^{k+1}\left(\frac{n!}{k!}\right) &= \sum_{k=1}^{n}(-1)^{k+1}(n!)\left(\frac{1}{k!}\right) \\
    &= n!\sum_{k=1}^{n}(-1)^{k+1}\left(\frac{1}{k!}\right) \\
    &= n!\left(\frac{1}{1!} - \frac{1}{2!} + \frac{1}{3!} - \dots \pm \frac{1}{n!}\right)
    \end{aligned}
    \]

    \medskip

    Then the number of derangements is the complement of the number of non-derangements, which is computed as follows:
    \[
    \begin{aligned}
    n! - n!\left(\frac{1}{1!} - \frac{1}{2!} + \frac{1}{3!} - \dots \pm \frac{1}{n!}\right) &= n! \cdot 1 - 
    n!\left(\frac{1}{1!} - \frac{1}{2!} + \frac{1}{3!} - \dots \pm \frac{1}{n!}\right) \\
    &= n!\left(1 - \left(\frac{1}{1!} - \frac{1}{2!} + \frac{1}{3!} - \dots \pm \frac{1}{n!}\right)\right) \\
    &= n!\left(1 - \frac{1}{1!} + \frac{1}{2!} - \frac{1}{3!} + \dots \pm \frac{1}{n!}\right)
    \end{aligned}
    \]

    \medskip

    Hence the number of non-derangements is
    \[
    n!\left(\frac{1}{1!} - \frac{1}{2!} + \frac{1}{3!} - \dots \pm \frac{1}{n!}\right)
    \]
    and the number of derangements is
    \[
    n!\left(1 - \frac{1}{1!} + \frac{1}{2!} - \frac{1}{3!} + \dots \pm \frac{1}{n!}\right)
    \]
\end{proof}

\begin{claim}
    As $n$ goes to infinity, the fraction of permutations that are derangements approaches $\frac{1}{e}$.
\end{claim}

\begin{proof}
    From Claim 6, the number of derangements can be written as
    \[
    n!\left(1 - \frac{1}{1!} + \frac{1}{2!} - \frac{1}{3!} + \dots \pm \frac{1}{n!}\right)
    \]
    thus the fraction of permutations that are derangements can be written as
    \[
    \frac{n!\left(1 - \frac{1}{1!} + \frac{1}{2!} - \frac{1}{3!} + \dots \pm \frac{1}{n!}\right)}{n!}
    = 1 - \frac{1}{1!} + \frac{1}{2!} - \frac{1}{3!} + \dots \pm \frac{1}{n!}
    \]

    \medskip

    Using the Taylor series for $e^x$,
    \[
    e^x = \sum_{k=0}^{\infty} \frac{x^k}{k!}
    \]
    we can substitute $x = -1$ to obtain
    \[
    e^{-1} = \sum_{k=0}^{\infty} \frac{(-1)^k}{k!} = 1 - \frac{1}{1!} + \frac{1}{2!} - \frac{1}{3!} + \dots
    \]
    
    \medskip
    
    Since
    \[
    1 - \frac{1}{1!} + \frac{1}{2!} - \frac{1}{3!} + \dots \pm \frac{1}{n!}
    \]
    is the $n$th partial sum of this series, it approaches $e^{-1} = \frac{1}{e}$ as $n$ goes to infinity.

    \medskip

    Hence the fraction of permutations that are derangements approaches $\frac{1}{e}$ as $n$ goes to infinity.
\end{proof}

\end{document}
